% Round-12 Autoformalism full-arm endpoints and saved baseline validation replay.
% Provenance and every underlying run are in docs/LLM_MEMORIZATION_EVIDENCE.md.
\subsection{Did the LLM Proposer Memorize the Model?}
\label{sec:experiment_memorization}

We examine dependence on familiar scientific names using the Dalla Man
$2\times2$ contrast: canonical versus perturbed dynamics, each presented with
named or obfuscated variables and descriptions. Within each semantic pair,
numerical trajectories, channel availability, and task burden are preserved.
Median validation free-rollout NMSE across two Autoformalism runs and three
GPT-5.6 Sol runs per cell is:

\begin{center}
\begin{tabular}{llrr}
\hline
Dynamics & Presentation & Autoformalism & GPT-5.6 Sol \\
\hline
Canonical & Named      & $0.354$ & $0.526$ \\
Canonical & Obfuscated & $0.257$ & $0.0310$ \\
Perturbed & Named      & $0.462$ & $0.0743$ \\
Perturbed & Obfuscated & $0.244$ & $4.74\times10^4$ \\
\hline
\end{tabular}
\end{center}

Autoformalism's median error does not increase under obfuscation for either
dynamics variant. GPT-5.6 Sol also improves under obfuscation on canonical
dynamics, so these results do not support a general failure when familiar names
are removed. Its large errors occur in the perturbed, obfuscated cell: two runs
reach NMSE $4.74\times10^4$ and $7.81\times10^4$, while the third reaches $0.100$.
Autoformalism's two errors in that cell are $0.245$ and $0.244$. Thus the observed
contrast concerns reliability under combined dynamics and presentation changes,
with substantial variation across GPT-5.6 Sol runs.

These results support Autoformalism's robustness to the tested changes without
relying on explicit physiological naming cues. They neither prove that its
proposer avoids memorization nor establish that GPT-5.6 Sol retrieved a known
model. Obfuscation does not remove every recognizable pattern, and differences
in search, numerical fitting, and model selection can also affect robustness.
The small, unequal run counts and reuse of validation for model selection make
this a descriptive development comparison.
